
\documentclass{article}
\usepackage{colortbl}
\usepackage{makecell}
\usepackage{multirow}
\usepackage{supertabular}

\begin{document}

\newcounter{utterance}

\centering \large Interaction Transcript for game `referencegame', experiment `number\_grids', episode 3 with NEW\_Reflect\_Agent.
\vspace{24pt}

{ \footnotesize  \setcounter{utterance}{1}
\setlength{\tabcolsep}{0pt}
\begin{supertabular}{c@{$\;$}|p{.15\linewidth}@{}p{.15\linewidth}p{.15\linewidth}p{.15\linewidth}p{.15\linewidth}p{.15\linewidth}}
    \# & $\;$A & \multicolumn{4}{c}{Game Master} & $\;\:$B\\
    \hline

    \theutterance \stepcounter{utterance}  
    & & \multicolumn{4}{p{0.6\linewidth}}{
        \cellcolor[rgb]{0.9,0.9,0.9}{
            \makecell[{{p{\linewidth}}}]{
                \texttt{\tiny{[P1$\langle$GM]}}
                \texttt{You are given three grids, where each of them is 5 by 5 in size.} \\
\texttt{Grids have empty cells marked with "▢" and filled cells marked with "X".} \\
\texttt{Your task is to generate a referring expression that best describes the target grid while distinguishing it from the two other distractor grids.} \\
\texttt{The first grid is the target grid, and the following two grids are the distractors.} \\
\\ 
\texttt{Target grid:} \\
\\ 
\texttt{□ X X X X} \\
\texttt{□ X □ □ X} \\
\texttt{□ X X X X} \\
\texttt{□ □ □ □ □} \\
\texttt{□ X X X X} \\
\\ 
\texttt{Distractor grid 1:} \\
\\ 
\texttt{□ X X X X} \\
\texttt{□ X □ □ X} \\
\texttt{□ X X X X} \\
\texttt{□ X □ □ X} \\
\texttt{□ X X X X} \\
\\ 
\texttt{Distractor grid 2:} \\
\\ 
\texttt{□ X X X X} \\
\texttt{□ X □ □ X} \\
\texttt{□ X □ □ X} \\
\texttt{□ X □ □ X} \\
\texttt{□ X X X X} \\
\\ 
\texttt{Instruction: Describe the target grid.} \\
\texttt{Generate the referring expression starting with the tag "Expression: " for the given target grid. Omit any other text.} \\
            }
        }
    }
    & \\ \\

    \theutterance \stepcounter{utterance}  
    & \multicolumn{4}{p{0.45\linewidth}}{
        \cellcolor[rgb]{0.95,0.95,0.95}{
            \makecell[{{p{\linewidth}}}]{
                \texttt{\tiny{[P1$|$internal]}}
                \texttt{Reflections from previous episodes:} \\
\texttt{{-}} \\
\texttt{ Failed. I described the target as having "full columns of Xs in columns 2{-}5," but that was true for all grids. The key difference was that the target has **every cell in the second column filled (all Xs)**, while Distractor 1 has an empty cell in row 4, column 2, and Distractor 2 has empty cells in rows 3{-}4, column 2. I should have focused specifically on the **second column being completely solid** to distinguish it. Next time, I will compare columns directly and pinpoint the exact column and its unique property.} \\
\texttt{{-}} \\
\texttt{ Expression: The grid where the leftmost column is all empty squares and the pattern of filled cells in rows 2 and 4 matches the pattern in rows 1, 3, and 5.} \\
\\ 
\texttt{Reflection: Failed. The expression was ambiguous; it did not clearly rule out Distractor 2, which also has a fully empty left column and similar row patterns. I should have specified that only rows 2 and 4 have a break in the filled cells (specifically, columns 3 and 4 are empty), unlike Distractor 2 where those rows are fully filled. Next time, I will explicitly note the exact positions of empty cells in the middle rows to distinguish the target uniquely.} \\
\texttt{{-}} \\
\texttt{ I succeeded. The target grid has a vertical column of X's in the first column except for the fourth row, which is empty. The distractors differ: distractor 1 has an X in the first column of row 4, and distractor 2 has the pattern shifted right. My expression highlighted the unique empty cell at (row 4, column 1). Next time, I would similarly focus on the single distinguishing empty cell in an otherwise filled column.} \\
            }
        }
    }
    & & \\ \\

    \theutterance \stepcounter{utterance}  
    & \multicolumn{4}{p{0.45\linewidth}}{
        \cellcolor[rgb]{0.95,0.95,0.95}{
            \makecell[{{p{\linewidth}}}]{
                \texttt{\tiny{[P1$|$internal]}}
                \texttt{{[}attempt 1 – valid{]}} \\
\texttt{Expression: The grid where the first column has Xs in all rows except row 4, which is empty, and rows 2 and 4 have empty cells in columns 3 and 4, while row 3 has all Xs in columns 2{-}5.} \\
            }
        }
    }
    & & \\ \\

    \theutterance \stepcounter{utterance}  
    & \multicolumn{4}{p{0.6\linewidth}}{
        \cellcolor[rgb]{0.9,0.9,0.9}{
            \makecell[{{p{\linewidth}}}]{
                \texttt{\tiny{[P1$\rangle$GM]}}
                \texttt{Expression: The grid where the first column has Xs in all rows except row 4, which is empty, and rows 2 and 4 have empty cells in columns 3 and 4, while row 3 has all Xs in columns 2{-}5.} \\
            }
        }
    }
    & & \\ \\

    \theutterance \stepcounter{utterance}  
    & & & \multicolumn{2}{p{0.3\linewidth}}{
        \cellcolor[rgb]{0.9,0.9,0.9}{
            \makecell[{{p{\linewidth}}}]{
                \texttt{\tiny{[GM$|$GM]}}
                \texttt{Expression: The grid where the first column has Xs in all rows except row 4, which is empty, and rows 2 and 4 have empty cells in columns 3 and 4, while row 3 has all Xs in columns 2{-}5.} \\
            }
        }
    }
    & & \\ \\

    \theutterance \stepcounter{utterance}  
    & & \multicolumn{4}{p{0.6\linewidth}}{
        \cellcolor[rgb]{0.9,0.9,0.9}{
            \makecell[{{p{\linewidth}}}]{
                \texttt{\tiny{[GM$\rangle$P2]}}
                \texttt{You are given three grids, where each of them is 5 by 5 in size.} \\
\texttt{Grids have empty cells marked with "▢" and filled cells marked with "X".} \\
\texttt{You are also given a referring expression that describes one of the given grids.} \\
\texttt{Your task is to select the grid that matches the given referring expression.} \\
\texttt{Generate only the number (in text) of the grid that matches the given expression by selecting first, second, or third.} \\
\\ 
\texttt{First:} \\
\\ 
\texttt{□ X X X X} \\
\texttt{□ X □ □ X} \\
\texttt{□ X X X X} \\
\texttt{□ □ □ □ □} \\
\texttt{□ X X X X} \\
\\ 
\texttt{Second:} \\
\\ 
\texttt{□ X X X X} \\
\texttt{□ X □ □ X} \\
\texttt{□ X X X X} \\
\texttt{□ X □ □ X} \\
\texttt{□ X X X X} \\
\\ 
\texttt{Third:} \\
\\ 
\texttt{□ X X X X} \\
\texttt{□ X □ □ X} \\
\texttt{□ X □ □ X} \\
\texttt{□ X □ □ X} \\
\texttt{□ X X X X} \\
\\ 
\texttt{Expression: The grid where the first column has Xs in all rows except row 4, which is empty, and rows 2 and 4 have empty cells in columns 3 and 4, while row 3 has all Xs in columns 2{-}5.} \\
\texttt{Question: Which grid does the expression refer to?} \\
\texttt{Start with the tag "Answer: ", followed by your selection. Omit any other text.} \\
            }
        }
    }
    & \\ \\

    \theutterance \stepcounter{utterance}  
    & & & \multicolumn{4}{p{0.45\linewidth}}{
        \cellcolor[rgb]{0.95,0.95,0.95}{
            \makecell[{{p{\linewidth}}}]{
                \texttt{\tiny{[P2$|$internal]}}
                \texttt{Reflections from previous episodes:} \\
\texttt{{-}} \\
\texttt{ Answer: Second} \\
\\ 
\texttt{Reflection: Succeeded. The expression described rows 1, 3, 5 fully filled, rows 2 and 4 with X only in first and last columns, and column 2 filled except row 4. Only the second grid matched all conditions. I carefully checked each row and column against the description. Next time, I will continue to verify each detail systematically to avoid misreading.} \\
\texttt{{-}} \\
\texttt{ Failed. I incorrectly selected the third grid, but the expression describes the first grid. I focused on "filled second column" but overlooked that the expression also says "the other two grids have at least one empty cell in that column." Checking all three:  } \\
\texttt{{-} First grid: second column is all X.  } \\
\texttt{{-} Second grid: second column has an empty in row 4.  } \\
\texttt{{-} Third grid: second column has empties in rows 3 and 4.  } \\
\texttt{So only first grid matches fully. Next time, verify the condition for *both* the target and the contrast with the other grids.} \\
\texttt{{-}} \\
\texttt{ Answer: Second} \\
\\ 
\texttt{Reflection: Succeeded. I carefully checked each grid against the expression's specific column and row conditions. The second grid matched all criteria: second column all Xs, first column Xs in rows 1,3,5 and empty in 2,4, and rows 2 and 4 had empty cells in columns 3 and 4. Next time, I will continue to verify each condition systematically to avoid oversight.} \\
            }
        }
    }
     \\ \\

    \theutterance \stepcounter{utterance}  
    & & & \multicolumn{4}{p{0.45\linewidth}}{
        \cellcolor[rgb]{0.95,0.95,0.95}{
            \makecell[{{p{\linewidth}}}]{
                \texttt{\tiny{[P2$|$internal]}}
                \texttt{{[}attempt 1 – valid{]}} \\
\texttt{Answer: Second} \\
            }
        }
    }
     \\ \\

    \theutterance \stepcounter{utterance}  
    & & & \multicolumn{4}{p{0.6\linewidth}}{
        \cellcolor[rgb]{0.9,0.9,0.9}{
            \makecell[{{p{\linewidth}}}]{
                \texttt{\tiny{[GM$\langle$P2]}}
                \texttt{Answer: Second} \\
            }
        }
    }
     \\ \\

    \theutterance \stepcounter{utterance}  
    & & & \multicolumn{2}{p{0.3\linewidth}}{
        \cellcolor[rgb]{0.9,0.9,0.9}{
            \makecell[{{p{\linewidth}}}]{
                \texttt{\tiny{[GM$|$GM]}}
                \texttt{second} \\
            }
        }
    }
    & & \\ \\

\end{supertabular}
}

\end{document}
